\subsection{2 段階 CCD 混合微分の精度解析：\texorpdfstring{$C^8$}{C\textasciicircum 8} 条件と劣化機構}
\label{app:ccd_mixed_precision}
% ═══════════════════════════════════════════════════════════════════════════════

§\ref{sec:ccd_mixed} の2段階 CCD 混合微分 $\phi_{xy}$ の精度を詳細に解析する．

Step 1 で得た $g = \partial f/\partial x$ の誤差は $E_g = O(h^6)$．
Step 2 で $g$ を $y$ 方向に微分するとき，誤差の上界は：
\[
  E_{\phi_{xy}} \leq \underbrace{O(h^6)}_{\text{$y$ 方向 CCD の打ち切り誤差}}
               + \underbrace{\left\|\pd{E_g}{y}\right\|}_{\text{$E_g$ の $y$ 方向変動}}
\]

\textbf{一般論では $\partial E_g/\partial y \sim O(h^6/h) = O(h^5)$ となる可能性がある．}
差分演算子は入力の誤差を $h^{-1}$ 倍程度増幅するため，$O(h^6)$ の誤差場を数値微分すると $O(h^5)$ に劣化するのが一般論である．

\textbf{ただし $f \in C^8$ であれば $O(h^6)$ が維持される：}
CCD の打ち切り誤差 $E_g(x,y) \propto h^6 f^{(7)}$ は $f$ の高次微分から生じる滑らかな場であり，
その $y$ 方向変動 $\partial E_g/\partial y \propto h^6 f^{(8)}$ は $f \in C^8$ のもとで $O(h^6)$ オーダーにとどまる．

\textbf{精度まとめ：}
\begin{itemize}
  \item $f \in C^8$ の場合：$E_{\phi_{xy}} = O(h^6)$ が維持される
  \item $f$ に不連続性や急峻な変化がある場合：精度が $O(h^5)$ 以下に劣化する可能性あり
\end{itemize}

\textbf{気液二相流の界面近傍への影響：}
気液二相流では，界面近傍で CLS 変数 $\psi$ が $O(\varepsilon)$ の厚さで急峻に変化するため，
$\psi$ の高次微分（$\psi^{(8)}$ 等）は $O(1/\varepsilon^7)$ のオーダーとなり，
実質的に $C^8$ 条件が界面近傍の格子点（界面から $3\varepsilon$ 以内）で成立しない．
したがって，曲率計算の混合微分 $\phi_{xy}$ の精度は界面近傍で $O(h^5)$ 以下に劣化しうる．
ただし，この劣化は界面近傍の $O(\varepsilon/h) \sim O(1)$ 個の格子点のみに影響し，
大域的な $L^2$ ノルム誤差への寄与は $O(h^{5+d/2})$（$d$：次元）程度に収まる
（比較：CSF モデル誤差 $O(\varepsilon^2) \approx O(h^2)$ が既に界面近傍精度を律速している）．
実務上は，格子収束テスト（第\ref{sec:verification}章）で曲率 $\kappa$ の収束次数を直接確認することが
この劣化の存在と程度を評価する最も確実な方法である．

実証的確認として，$f(x,y) = \sin(\pi x)\cos(\pi y)$（解析解 $f_{xy} = -\pi^2\cos(\pi x)\sin(\pi y)$）で格子収束テストを行い，$N$ を倍にするたびに誤差が $1/64$ 倍（6次収束）になるかを確認すること（第\ref{sec:verification}章参照）．

% ═══════════════════════════════════════════════════════════════════════════════
